% ============================================================
%  cap02.tex  --  Capítulo 2: Séries Finitas -- Capitalização Periódica
%  Baseado em: Notas de Aula 1 (8 páginas)
%  Referência: Bueno, Rangel & Santos, Cap. 2
% ============================================================

\chapter{Séries Finitas -- Capitalização Periódica}

\section{Visão Geral}

Uma \textbf{série de pagamentos} é uma sequência de fluxos de caixa $\{R_1, R_2, \ldots, R_n\}$ distribuída ao longo do tempo. Dependendo do comportamento desses valores, a série pode ser:

\medskip
\begin{center}
\begin{tikzpicture}[>=Stealth, font=\small,
  box/.style={draw, rounded corners=3pt, fill=blue!15,
              text width=2.5cm, align=center, font=\small\bfseries, inner sep=4pt},
  sbox/.style={draw, rounded corners=3pt, fill=gray!10,
               text width=2.3cm, align=center, font=\footnotesize, inner sep=3pt}
]
  % Column 1: tipos de série
  \node[sbox] (U)  at (0, 1.0)  {Uniforme};
  \node[sbox] (PA) at (0, 0.0)  {em PA};
  \node[sbox] (PG) at (0,-1.0)  {em PG};

  % Column 2: tipo de fluxo
  \node[sbox] (D)  at (4.0, 0.5)  {Discreta};
  \node[sbox] (Co) at (4.0,-0.5)  {Cont\'{i}nua};

  % Column 3: capitalização
  \node[sbox] (Per) at (8.0, 0.5)  {Peri\'{o}dica};
  \node[sbox] (Ins) at (8.0,-0.5)  {Instant\^{a}nea};

  % Arrows tipo -> fluxo
  \draw[->, thick, color=blue!60!black] (U.east)  -- (D.west);
  \draw[->, thick, color=blue!60!black] (PA.east) -- (D.west);
  \draw[->, thick, color=blue!60!black] (PG.east) -- (D.west);
  \draw[->, thick, color=blue!60!black] (U.east)  to[bend right=10] (Co.west);

  % Arrows fluxo -> capitalização
  \draw[->, thick, color=blue!60!black] (D.east)  -- (Per.west);
  \draw[->, thick, color=blue!60!black] (Co.east) -- (Ins.west);
\end{tikzpicture}
\end{center}

\medskip
Este capítulo trata de séries \textbf{discretas} sob \textbf{capitalização periódica} (juros compostos).

% ─────────────────────────────────────────────────────────────────────────────
\section{Soma da Progressão Geométrica}

As fórmulas de séries de pagamentos derivam diretamente da soma de uma Progressão Geométrica (PG).

\begin{definicao}
Uma \textbf{Progressão Geométrica} (PG) é uma sequência em que a razão entre termos consecutivos é constante e igual a $q$:
\[
  q = \frac{a_k}{a_{k-1}}, \qquad A = \{a_1,\; a_1 q,\; a_1 q^2,\; \ldots,\; a_1 q^{n-1}\}.
\]
\end{definicao}

\subsection*{Derivação da fórmula}

Seja $S$ a soma dos $n$ termos:
\begin{align}
  S    &= a_1 + a_1 q + a_1 q^2 + \cdots + a_1 q^{n-1}  \tag{I}  \\
  qS   &= \phantom{a_1 + {}} a_1 q + a_1 q^2 + \cdots + a_1 q^{n-1} + a_1 q^n  \tag{II}
\end{align}
Subtraindo (I) de (II):
\[
  qS - S = a_1 q^n - a_1
  \quad\Rightarrow\quad
  S(q-1) = a_1(q^n - 1).
\]

\begin{resultado}
\[
  \boxed{S = a_1 \left[\frac{q^n - 1}{q - 1}\right]}
  \qquad (q \neq 1).
\]
Para $q > 1$: PG \textbf{crescente}; \quad para $q < 1$: PG \textbf{decrescente}.
\end{resultado}

% ─────────────────────────────────────────────────────────────────────────────
\section{Séries Uniformes Vencidas (Postecipadas)}

\subsection{Diagrama de fluxo de caixa}

Na série uniforme vencida (\textit{ordinary annuity}), os pagamentos $R$ ocorrem ao \emph{final} de cada período:

\begin{center}
\begin{tikzpicture}[>=Stealth, font=\small, scale=0.95]
  \draw[timeline axis] (-0.5,0) -- (8.5,0) node[right] {$t$};
  \foreach \x/\lab in {0/0, 1/1, 2/2, 3/3, 7/\ensuremath{n}}{
    \draw[tick] (\x,0.12)--(\x,-0.12);
    \node[below=4pt] at (\x,0) {\lab};
  }
  \node at (5,-0.45) {$\cdots$};
  % Pagamentos
  \foreach \x in {1,2,3,7}{
    \draw[arrow down] (\x,-0.65)--(\x,-1.4);
    \node[cash label, below=2pt] at (\x,-1.4) {$R$};
  }
  % P (valor presente) -- seta para cima à esquerda
  \draw[arrow up] (0,0)--(0,1.3);
  \node[cash label, above=2pt] at (0,1.3) {$P$};
  % S (valor futuro) -- seta para baixo à direita
  \draw[arrow up, dashed, color=verde] (7,0)--(7,2.0);
  \node[cash label, above=2pt, color=verde] at (7,2.0) {$S$};
\end{tikzpicture}
\end{center}

\subsection{Fórmula do valor futuro (montante)}

Levando cada depósito $R$ ao instante $n$ (valor futuro):
\[
  S = R + R(1+i) + R(1+i)^2 + \cdots + R(1+i)^{n-1}.
\]
Esta é uma PG com $a_1 = R$, razão $q = (1+i)$ e $n$ termos. Aplicando a fórmula da soma da PG:
\begin{resultado}
\[
  \boxed{S = R\left[\frac{(1+i)^n - 1}{i}\right]}
\]
O fator $\dfrac{(1+i)^n-1}{i}$ é denominado \textbf{Fator de Acumulação de Capital (FAC)}.
\end{resultado}

\subsection{Fórmula do valor presente}

Trazendo $S$ ao instante zero:
\begin{resultado}
\[
  \boxed{P = R\left[\frac{(1+i)^n - 1}{i}\right] \cdot \frac{1}{(1+i)^n}
           = \frac{R}{(1+i)} + \frac{R}{(1+i)^2} + \cdots + \frac{R}{(1+i)^n}}
\]
\end{resultado}

Isolando $R$ (prestação de um financiamento de principal $P$):
\[
  R = P\left[\frac{(1+i)^n \cdot i}{(1+i)^n - 1}\right].
\]

\subsection*{Share de juros}

O \textbf{share de juros} $j/P$ mede a fração do valor presente que corresponde a juros:
\[
  \frac{J}{P} = \frac{S - P}{P} = (1+i)^n - 1.
\]

\begin{exemplo}{Poupança mensal -- acumulação}
\label{ex:poupança}
Depósitos mensais de R\$\,150 durante 20 meses à taxa de $2\%$ a.m.
\[
  S = 150\left[\frac{(1{,}02)^{20}-1}{0{,}02}\right]
    = 150 \times 24{,}2974 \approx \textbf{R\$\,3.644{,}61.}
\]
\end{exemplo}

\begin{exemplo}{Financiamento -- cálculo da prestação}
\label{ex:financiamento}
Financiamento de R\$\,25.000 em 24 prestações mensais iguais à taxa de $3\%$ a.m.
\[
  R = 25{.}000\left[\frac{(1{,}03)^{24} \times 0{,}03}{(1{,}03)^{24}-1}\right]
    \approx \textbf{R\$\,1.476{,}19.}
\]
\end{exemplo}

\begin{moderno}[title={Financiamento imobiliário em 2025}]
Considere um imóvel financiado pelo sistema PRICE com:
\begin{itemize}[leftmargin=1.8em, nosep]
  \item Principal: R\$\,300.000
  \item Prazo: 360 meses (30 anos)
  \item Taxa efetiva: $0{,}79\%$ a.m. (equivalente à TR + $9{,}89\%$ a.a. -- referência Caixa, 2025)
\end{itemize}
\[
  R = 300{.}000 \times \frac{(1{,}0079)^{360}\times 0{,}0079}{(1{,}0079)^{360}-1}
    \approx \textbf{R\$\,2.547{,}48/mês.}
\]
O total pago ao final ($360 \times 2.547{,}48 = \text{R\$\,}917.092{,}80$) é aproximadamente
3 vezes o principal.
\end{moderno}

% ─────────────────────────────────────────────────────────────────────────────
\section{Série Antecipada}

Na série antecipada, os pagamentos ocorrem no \emph{início} de cada período (em vez do final). A única diferença em relação à série vencida é o avanço de um período no fluxo:

\begin{center}
\begin{tikzpicture}[>=Stealth, font=\small, scale=0.95]
  \draw[timeline axis] (-0.5,0) -- (8.5,0) node[right] {$t$};
  \foreach \x/\lab in {0/0, 1/1, 2/2, 3/3, 7/\ensuremath{n}}{
    \draw[tick] (\x,0.12)--(\x,-0.12);
    \node[below=4pt] at (\x,0) {\lab};
  }
  \node at (5,-0.35) {$\cdots$};
  \foreach \x in {0,1,2,3}{
    \draw[arrow down] (\x,-0.65)--(\x,-1.45);
    \node[cash label, below=2pt] at (\x,-1.45) {$R$};
  }
  \draw[arrow up, dashed, color=verde] (7,0)--(7,2.0);
  \node[cash label, above=2pt, color=verde] at (7,2.0) {$S$};
\end{tikzpicture}
\end{center}

Cada depósito fica aplicado por um período a mais, portanto:
\begin{resultado}
\[
  S_{\text{ant}} = R\left[\frac{(1+i)^n-1}{i}\right](1+i), \qquad
  P_{\text{ant}} = R\left[\frac{(1+i)^n-1}{i}\right]\frac{1+i}{(1+i)^n}.
\]
\end{resultado}
A prestação de um financiamento antecipado é:
\[
  R_{\text{ant}} = P\left[\frac{(1+i)^n \cdot i}{(1+i)^n - 1}\right]\frac{1}{1+i}.
\]

% ─────────────────────────────────────────────────────────────────────────────
\section{Séries com Pagamentos Variáveis}

Quando os pagamentos $R_t$ variam ao longo do tempo, não existe fórmula fechada genérica: calcula-se diretamente:
\[
  S = \sum_{t=1}^{n} R_t\,(1+i)^{n-t},
  \qquad
  P = \sum_{t=1}^{n} \frac{R_t}{(1+i)^t}.
\]

\begin{exemplo}{5 pagamentos de R\$\,1.000 + 5 de R\$\,2.000 a $2\%$ a.m.}
\[
  P = \frac{1.000}{1{,}02} + \frac{1.000}{(1{,}02)^2} + \cdots
      + \frac{2.000}{(1{,}02)^9} + \frac{2.000}{(1{,}02)^{10}}
    \approx \textbf{R\$\,13.251{,}71.}
\]
\end{exemplo}

\begin{exemplo}{Empréstimo de R\$\,1.200 -- taxa implícita}
Crédito de R\$\,1.200 a vista; pagamento de $6 \times \text{R\$\,100} + 6 \times \text{R\$\,150}$.
Resolva para $i$:
\[
  1.200 = \frac{100}{(1+i)} + \frac{100}{(1+i)^2} + \cdots
          + \frac{150}{(1+i)^{11}} + \frac{150}{(1+i)^{12}}
  \quad\Rightarrow\quad i \approx 3{,}28\%\text{ a.m.}
\]
(Resolução numérica -- ver Capítulo 7, Método de Newton--Raphson.)
\end{exemplo}

% ─────────────────────────────────────────────────────────────────────────────
\section{Séries em Progressão Aritmética (Gradiente)}

\begin{definicao}
Uma série em \textbf{Progressão Aritmética} (PA) possui pagamentos que crescem (ou decrescem) de um valor constante $R$ em cada período. O primeiro pagamento é igual à \textbf{razão} $R$:
\[
  \{R,\; 2R,\; 3R,\; \ldots,\; nR\}.
\]
\end{definicao}

\subsection{Gradiente crescente}

\begin{center}
\begin{tikzpicture}[>=Stealth, font=\small, scale=0.95]
  \draw[timeline axis] (-0.5,0) -- (8.5,0) node[right] {$t$};
  \foreach \x/\lab in {0/0, 1/1, 2/2, 3/3, 7/\ensuremath{n}}{
    \draw[tick] (\x,0.12)--(\x,-0.12);
    \node[below=4pt] at (\x,0) {\lab};
  }
  \node at (5.5,-0.35) {$\cdots$};
  \draw[arrow down] (1,-0.65)--(1,-1);  \node[cash label,below=2pt] at(1,-0.9){$R$};
  \draw[arrow down] (2,-0.65)--(2,-1.5);  \node[cash label,below=2pt] at(2,-1.4){$2R$};
  \draw[arrow down] (3,-0.65)--(3,-2);  \node[cash label,below=2pt] at(3,-1.9){$3R$};
  \draw[arrow down] (7,-0.65)--(7,-3);  \node[cash label,below=2pt] at(7,-2.9){$nR$};
  \draw[arrow up]   (0,0)--(0,1.3);   \node[cash label,above=2pt] at(0,1.3){$P$};
\end{tikzpicture}
\end{center}

\subsection*{Derivação do valor futuro}

Tomando o valor futuro da série $S = R(1+i)^{n-1} + 2R(1+i)^{n-2} + \cdots + nR$ e aplicando o mesmo truque da soma de PG (multiplicar por $(1+i)$ e subtrair), chega-se a:

\begin{resultado}
\[
  S = \frac{R}{i}\left\{(1+i)\left[\frac{(1+i)^n-1}{i}\right] - n\right\}.
\]
\end{resultado}

O valor presente correspondente é:
\begin{resultado}
\[
  P = \frac{R}{i}\left\{(1+i)\left[\frac{(1+i)^n-1}{i}\right] - n\right\}
      \cdot\frac{1}{(1+i)^n}.
\]
\end{resultado}

\begin{exemplo}{Depósitos crescentes em PA a $2{,}5\%$ a.m.}
Depósitos de R\$\,100 no 1º mês, R\$\,200 no 2º, \ldots, R\$\,3.600 no 36º mês ($R = 100$, $n = 36$, $i = 0{,}025$):
\[
  S = \frac{100}{0{,}025}\left\{1{,}025\left[\frac{(1{,}025)^{36}-1}{0{,}025}\right]-36\right\}
    \approx \textbf{R\$\,90.935{,}79.}
\]
\end{exemplo}

\subsection{Gradiente decrescente}

Na série em PA \textbf{decrescente}, os pagamentos são $\{nR, (n-1)R, \ldots, 2R, R\}$:

\begin{center}
\begin{tikzpicture}[>=Stealth, font=\small, scale=0.95]
  \draw[timeline axis] (-0.5,0) -- (8.5,0) node[right] {$t$};
  \foreach \x/\lab in {0/0, 1/1, 2/2, 3/3, 7/\ensuremath{n}}{
    \draw[tick] (\x,0.12)--(\x,-0.12);
    \node[below=4pt] at (\x,0) {\lab};
  }
  \node at (5.5,-0.35) {$\cdots$};
  \draw[arrow down] (1,-0.65)--(1,-2.8);  \node[cash label,below=2pt] at(1,-2.9){$nR$};
  \draw[arrow down] (2,-0.65)--(2,-2.3);  \node[cash label,below=2pt] at(2,-2.4){$(n{-}1)R$};
  \draw[arrow down] (3,-0.65)--(3,-1.8);  \node[cash label,below=2pt] at(3,-1.9){$(n{-}2)R$};
  \draw[arrow down] (7,-0.65)--(7,-1);  \node[cash label,below=2pt] at(7,-0.9){$R$};
  \draw[arrow up]   (0,0)--(0,1.3);   \node[cash label,above=2pt] at(0,1.3){$P$};
\end{tikzpicture}
\end{center}

Usando o mesmo procedimento:
\begin{resultado}
\[
  S = \frac{R}{i}\left\{n(1+i)^n - \left[\frac{(1+i)^n-1}{i}\right]\right\},
  \qquad
  P = \frac{R}{i}\left\{n(1+i)^n - \left[\frac{(1+i)^n-1}{i}\right]\right\}
      \frac{1}{(1+i)^n}.
\]
\end{resultado}

% ─────────────────────────────────────────────────────────────────────────────
\section{Série com Valor Inicial Diferente da Razão}

Seja uma série com primeiro pagamento $F$ e razão $R$ (i.e., pagamentos $F{+}R$, $F{+}2R$, \ldots, $F{+}nR$):

\begin{center}
\begin{tikzpicture}[>=Stealth, font=\small, scale=0.95]
  \draw[timeline axis] (-0.5,0) -- (8.5,0) node[right] {$t$};
  \foreach \x/\lab in {0/0, 1/1, 2/2, 3/3, 4/4, 8/\ensuremath{n}}{
    \draw[tick] (\x,0.12)--(\x,-0.12);
    \node[below=4pt] at (\x,0) {\lab};
  }
  \node at (6.5,-0.35) {$\cdots$};
  \draw[arrow down] (1,-0.65)--(1,-1.10);  \node[cash label,below=2pt] at(1,-1.1){$F{+}R$};
  \draw[arrow down] (2,-0.65)--(2,-1.6);  \node[cash label,below=2pt] at(2,-1.6){$F{+}2R$};
  \draw[arrow down] (3,-0.65)--(3,-2.1);  \node[cash label,below=2pt] at(3,-2.1){$F{+}3R$};
  \draw[arrow down] (8,-0.65)--(8,-2.9);  \node[cash label,below=2pt] at(8,-2.9){$F{+}nR$};
  \draw[arrow up]   (0,0)--(0,1.3);   \node[cash label,above=2pt] at(0,1.3){$P$};
\end{tikzpicture}
\end{center}

A estratégia é decompor em uma série gradiente (razão $R$) mais uma série uniforme (valor fixo $F$):
\begin{resultado}
\[
  S = \underbrace{\frac{R}{i}\left\{(1+i)\left[\frac{(1+i)^n-1}{i}\right]-n\right\}}_{S_{\text{grad}}}
    + \underbrace{F\left[\frac{(1+i)^n-1}{i}\right]}_{S_{\text{unif}}}
    \quad\text{(postecipada)}.
\]
\end{resultado}

\begin{exemplo}{Série com $F{+}R = 500$, $R = 50$, $F = 400$, $n = 35$, $i = 2\%$ a.m.}
\[
  S_{\text{grad}}
    = \frac{50}{0{,}02}\left\{1{,}02\left[\frac{(1{,}02)^{35}-1}{0{,}02}\right]-35\right\}
    \approx 17{.}933{,}57.
\]
\[
  S_{\text{unif}} = 400\left[\frac{(1{,}02)^{35}-1}{0{,}02}\right]
                 \approx 13{.}029{,}39.
\]
\[
  S_{\text{total}} = 17{.}933{,}57 + 13{.}029{,}39 \approx \textbf{R\$\,30.962{,}96.}
\]
\end{exemplo}

% ─────────────────────────────────────────────────────────────────────────────
\section{Séries em Progressão Geométrica}

\begin{definicao}
Em uma série em \textbf{PG}, cada pagamento cresce (ou decresce) pelo fator constante $(1+\alpha)$:
\[
  \{R,\; R(1+\alpha),\; R(1+\alpha)^2,\; \ldots,\; R(1+\alpha)^{n-1}\}.
\]
Para $\alpha > 0$: série \textbf{crescente}; $\alpha < 0$: \textbf{decrescente}; $\alpha = 0$: série uniforme.
\end{definicao}

\begin{center}
\begin{tikzpicture}[>=Stealth, font=\small, scale=0.95]
  \draw[timeline axis] (-0.5,0) -- (8.5,0) node[right] {$t$};
  \foreach \x/\lab in {0/0, 1/1, 2/2, 3/3, 8/\ensuremath{n}}{
    \draw[tick] (\x,0.12)--(\x,-0.12);
    \node[below=4pt] at (\x,0) {\lab};
  }
  \node at (6,-0.35) {$\cdots$};
  \draw[arrow down] (1,-0.65)--(1,-1);
    \node[cash label,below=2pt] at(1,-1.0){$R$};
  \draw[arrow down] (2,-0.65)--(2,-1.55);
    \node[cash label,below=2pt] at(2,-1.5){$R(1{+}\alpha)$};
  \draw[arrow down] (3,-0.65)--(3,-2.15);
    \node[cash label,below=2pt] at(3,-2.0){$R(1{+}\alpha)^2$};
  \draw[arrow down] (8,-0.65)--(8,-3);
    \node[cash label,below=2pt] at(8,-2.9){$R(1{+}\alpha)^{n-1}$};
  \draw[arrow up]   (0,0)--(0,1.3);
    \node[cash label,above=2pt] at(0,1.3){$P$};
\end{tikzpicture}
\end{center}

O valor futuro é uma PG com razão $\delta = \dfrac{1+\alpha}{1+i}$:
\[
  S = R\,(1+i)^{n-1}\left[\frac{\left(\frac{1+\alpha}{1+i}\right)^n - 1}{\frac{1+\alpha}{1+i}-1}\right].
\]
Simplificando:
\begin{resultado}
\[
  S = R\left[\frac{(1+i)^n - (1+\alpha)^n}{i - \alpha}\right], \qquad i \neq \alpha.
\]
\end{resultado}

O valor presente correspondente é:
\[
  P = \frac{R}{i-\alpha}\left[1 - \left(\frac{1+\alpha}{1+i}\right)^n\right].
\]

\begin{nota}
\textbf{Caso especial $i = \alpha$:} A expressão acima é indeterminada. Neste caso, todos os fluxos trazidos a valor futuro são iguais, portanto:
\[
  S = n\,R\,(1+i)^{n-1}.
\]
\end{nota}

\begin{moderno}[title={Salários crescentes e poupança para aposentadoria}]
Um trabalhador deposita mensalmente o equivalente a $5\%$ do salário, que cresce $0{,}5\%$ ao mês ($\alpha = 0{,}005$). O salário inicial é R\$\,4.000, portanto $R_1 = 200$. Taxa de rendimento: $0{,}8\%$ a.m. ($i = 0{,}008$). Após 30 anos ($n = 360$ meses):
\[
  S = 200 \times \frac{(1{,}008)^{360} - (1{,}005)^{360}}{0{,}008 - 0{,}005}
    \approx \textbf{R\$\,2.874.300.}
\]
A diferença entre usar taxa de crescimento salarial vs.\ depósito fixo é expressiva.
\end{moderno}
